\section{MPO tensors}

\begin{definition}[MPO tensor]\label{def:mpo_tensor}
    \lean{MPOTensor}
    \leanok
    An \emph{MPO tensor} is a family of matrices
    $\{M^{ij}\}_{i,j=0}^{d-1}$ with $M^{ij} \in \MN{D}$, indexed by a ket
    index $i$ and a bra index $j$. Diagrammatically,
    \begin{center}
        \tenkzkernel
        % Formula: (M^{ij})_{\alpha\beta}, with M^{ij}\in\operatorname{Mat}_D.
        % Ink/index correspondence: the west and east legs carry the free
        % virtual indices \alpha and \beta; the upper and lower legs carry the
        % free physical indices i and j.  Nothing is contracted, so the
        % boundary signature (virtual-west,virtual-east,physical-up,physical-down)
        % is (1,1,1,1).
        \begin{tenkz}[rows={wire}, cols=1, west=open, east=open]
            \tn[skin=mpo, ports={90:physical:$i$, 270:physical:$j$}]{M}
        \end{tenkz}
    \end{center}
\end{definition}

\begin{definition}[Doubled-index MPS view]\label{def:mpo_to_mps}
    \lean{MPOTensor.toMPSTensor}
    \leanok
    \uses{def:mpo_tensor}
    By combining the pair $(i,j)$ into a single physical index in
    $\{0,\ldots,d^2-1\}$, one obtains an MPS tensor with physical dimension
    $d^2$ and bond dimension $D$, with
    $M^{\mathrm{MPS}}_{(i,j)}=M^{ij}$, identifying
    $\{0,\ldots,d{-}1\}\times\{0,\ldots,d{-}1\}$ with
    $\{0,\ldots,d^2{-}1\}$ via the product encoding.
\end{definition}

\begin{definition}[Word evaluation]\label{def:mpo_eval_word}
    \lean{MPOTensor.evalWord}
    \leanok
    \uses{def:mpo_tensor}
    For words $u=(i_1,\ldots,i_N)$ and $v=(j_1,\ldots,j_N)$ in
    $\{0,\ldots,d{-}1\}^N$, the \emph{word evaluation} is
    \begin{align}
        M^{u,v}
        &:=M^{i_1j_1}M^{i_2j_2}\cdots M^{i_Nj_N}\in\MN{D}.
        \label{eq:mpdo_word_eval}
    \end{align}
    The empty pair of words evaluates to the identity,
    $M^{\varnothing,\varnothing}=\Id_D$, and word pairs of different
    lengths evaluate to $0$.
\end{definition}

\begin{lemma}[Word-evaluation factorization]\label{lem:mpo_eval_word_append}
    \lean{MPOTensor.evalWord_append}
    \leanok
    \uses{def:mpo_eval_word}
    For words $u_1,v_1$ of equal length and arbitrary words $u_2,v_2$,
    word evaluation factors as
    $M^{u_1u_2,\,v_1v_2}=M^{u_1,v_1}M^{u_2,v_2}$.
\end{lemma}

\begin{proof}\leanok\uses{def:mpo_eval_word}
    We argue by induction on the common length of $u_1$ and $v_1$. For the
    empty prefix, $M^{u_2,v_2}=\Id_D M^{u_2,v_2}$. If $u_1=iu$ and
    $v_1=jv$, then
    \begin{align}
        M^{iuu_2,\,jvv_2}
        &=M^{ij}M^{uu_2,\,vv_2}
          =M^{ij}M^{u,v}M^{u_2,v_2}
          =M^{iu,jv}M^{u_2,v_2},
        \notag
    \end{align}
    using associativity of matrix multiplication.
\end{proof}

\begin{lemma}[Trace cyclicity of the closed MPO word]\label{lem:mpo_trace_evalWord_cyclic}
    \lean{MPOTensor.trace_evalWord_cons_eq_append}
    \leanok
    \uses{def:mpo_eval_word}
    Let $a,b\in\{0,\ldots,d{-}1\}$, and let $u,v$ be words of equal length.
    Writing $au,bv$ for adding the letters at the beginning and $ua,vb$ for
    adding them at the end, one has
    $\tr\!\left(M^{au,\,bv}\right)=\tr\!\left(M^{ua,\,vb}\right)$.
\end{lemma}

\begin{proof}\leanok\uses{def:mpo_eval_word, lem:mpo_eval_word_append}
    Lemma~\ref{lem:mpo_eval_word_append} gives
    $M^{ua,\,vb}=M^{u,v}M^{ab}$ and
    $M^{au,\,bv}=M^{ab}M^{u,v}$. Hence, by cyclicity of the trace,
    \begin{align}
        \tr\!\left(M^{ua,\,vb}\right)
        &=\tr\!\left(M^{u,v}M^{ab}\right)
          =\tr\!\left(M^{ab}M^{u,v}\right)
          =\tr\!\left(M^{au,\,bv}\right).
        \notag
    \end{align}
\end{proof}

\begin{definition}[MPO operator family]\label{def:mpo_family}
    \lean{MPOTensor.mpo}
    \leanok
    \uses{def:mpo_eval_word}
    The operator generated by $M$ on $N$ sites is the matrix
    $\rho^{(N)}(M)$ indexed by configurations
    $\sigma,\tau\in\{0,\ldots,d{-}1\}^N$ and given by
    \begin{align}
        \rho^{(N)}(M)_{\sigma,\tau}
        &=\tr(M^{\sigma,\tau}).
        \label{eq:mpdo_family}
    \end{align}
\end{definition}

Diagrammatically,
\begin{center}
    \tenkzkernel
    % Formula: \rho^{(N)}(M)_{\sigma,\tau}
    % = \tr(M^{\sigma_1\tau_1}\cdots
    % M^{\sigma_N\tau_N})
    % = \sum_{\alpha_1,\ldots,\alpha_N}
    % M^{\sigma_1\tau_1}_{\alpha_1\alpha_2}\cdots
    % M^{\sigma_N\tau_N}_{\alpha_N\alpha_1}.
    % Ink/index correspondence: every horizontal virtual index is contracted,
    % including the return bond \alpha_{N+1}=\alpha_1.  The upper legs carry
    % the free indices \sigma_1,\ldots,\sigma_N and the lower legs carry
    % \tau_1,\ldots,\tau_N, hence the semantic boundary signature is
    % (0,0,N,N).  The \tndots atom abbreviates the omitted cells, so the drawn
    % and logged physical arities are necessarily schematic rather than N.
    \begin{tenkz}[rows={wire}, cols=4, boundary=periodic]
        \tn[skin=mpo, ports={90:physical:$\sigma_1$, 270:physical:$\tau_1$}]{M} &
        \tn[skin=mpo, ports={90:physical:$\sigma_2$, 270:physical:$\tau_2$}]{M} &
        \tn[skin=dots]{} &
        \tn[skin=mpo, ports={90:physical:$\sigma_N$, 270:physical:$\tau_N$}]{M}
    \end{tenkz}
\end{center}

\begin{definition}[Cyclic-shift reindexing]\label{def:mpo_rotate_config}
    \lean{MPOTensor.rotateConfig}
    \leanok
    The \emph{cyclic shift} of configurations on $N$ sites is the bijection that
    relabels each site by one position, $\sigma\mapsto\sigma\circ r$, where
    $r$ is the cyclic rotation $i\mapsto i+1$ of $\{0,\ldots,N-1\}$.
\end{definition}

\begin{theorem}[Translation invariance of the periodic MPDO]
    \label{thm:mpo_translation_invariant}
    \lean{MPOTensor.mpo_submatrix_rotateConfig}
    \leanok
    \uses{def:mpo_family, def:mpo_rotate_config}
    The operator $\rho^{(N)}(M)$ is invariant under the simultaneous cyclic shift
    of its bra and ket configurations:
    $\rho^{(N)}(M)_{r\sigma,\,r\tau}=\rho^{(N)}(M)_{\sigma,\tau}$.
\end{theorem}

\begin{proof}\leanok
    The entry is a closed bond trace,
    $\rho^{(N)}(M)_{\sigma,\tau}
    =\tr(M^{\sigma_1\tau_1}\cdots M^{\sigma_N\tau_N})$.
    Shifting both configurations by $r$ rotates the matrix product by one factor,
    and the trace is invariant under that cyclic rotation:
    \begin{align}
        \tr\bigl(        M^{\sigma_2\tau_2}\cdots
            M^{\sigma_N\tau_N}M^{\sigma_1\tau_1}
        \bigr)
        &=
        \tr\bigl(        M^{\sigma_1\tau_1}\cdots M^{\sigma_N\tau_N}
        \bigr),
        \notag
    \end{align}
    so $\rho^{(N)}(M)_{r\sigma,\,r\tau}
    =\rho^{(N)}(M)_{\sigma,\tau}$.
\end{proof}

\begin{definition}[Hermitian MPO tensor]\label{def:mpo_hermitian}
    \lean{MPOTensor.IsHermitian}
    \leanok
    \uses{def:mpo_tensor}
    An MPO tensor is \emph{Hermitian} if
    $M^{ij}=(M^{ji})^\dagger$ for all
    $i,j\in\{0,\ldots,d{-}1\}$.
\end{definition}

\begin{definition}[Adjoint MPO tensor]\label{def:mpo_adjoint_tensor}
    \lean{MPOTensor.adjointTensor}
    \lean{MPOTensor.adjointTensor_eq_iff_isHermitian}
    \leanok
    \uses{def:mpo_tensor, def:mpo_hermitian}
    The \emph{adjoint tensor} $M^\dagger$ of an MPO tensor $M$ exchanges the
    ket and bra physical indices and conjugate-transposes each virtual
    matrix: $(M^\dagger)^{ij}=(M^{ji})^\dagger$.
    A tensor equals its adjoint tensor, $M^\dagger=M$, precisely when it
    is Hermitian.
\end{definition}

\begin{theorem}[Reflection identity for the adjoint tensor]
    \label{thm:mpo_adjoint_reflection}
    \lean{MPOTensor.evalWord_adjointTensor}
    \lean{MPOTensor.mpo_adjointTensor}
    \lean{MPOTensor.IsMPDO.mpo_adjointTensor_eq}
    \lean{MPOTensor.IsMPDO.adjointTensor}
    \leanok
    \uses{def:mpo_adjoint_tensor, def:mpo_family, def:mpdo}
    The adjoint tensor generates the adjoint operator family up to spatial
    reflection: entrywise,
    \begin{align}
        \rho^{(N)}(M^\dagger)_{\sigma,\tau}
        &=
        \overline{
            \rho^{(N)}(M)_{\tau\circ\mathrm{rev},
            \sigma\circ\mathrm{rev}}
        },
        \label{eq:mpdo_coeff_adjoint_reflection}
    \end{align}
    where $\mathrm{rev}$ reverses the site order. For an MPDO and $N\geq1$,
    the conjugation disappears:
    \begin{align}
        \rho^{(N)}(M^\dagger)_{\sigma,\tau}
        &=
        \rho^{(N)}(M)_{\sigma\circ\mathrm{rev},
        \tau\circ\mathrm{rev}},
        \label{eq:mpdo_adjoint_mpdo}
    \end{align}
    and the adjoint tensor is again an MPDO.
\end{theorem}

\begin{proof}
    \leanok
    At the level of word evaluations, the conjugate transpose of a product
    reverses the factors, so
    \begin{align}
        (M^\dagger)^{\sigma,\tau}
        &=
        \bigl(        M^{\tau\circ\mathrm{rev},
            \sigma\circ\mathrm{rev}}
        \bigr)^\dagger.
        \notag
    \end{align}
    Taking traces gives the entrywise identity
    \begin{align}
        \rho^{(N)}(M^\dagger)_{\sigma,\tau}
        &=
        \tr\!\bigl(        (M^{\tau\circ\mathrm{rev},
            \sigma\circ\mathrm{rev}})^\dagger
        \bigr)
        =
        \overline{
            \rho^{(N)}(M)_{\tau\circ\mathrm{rev},
            \sigma\circ\mathrm{rev}}
        }.
        \notag
    \end{align}
    For an MPDO, $\rho^{(N)}(M)$ is Hermitian, so
    \begin{align}
        \overline{
            \rho^{(N)}(M)_{\tau\circ\mathrm{rev},
            \sigma\circ\mathrm{rev}}
        }
        &=
        \rho^{(N)}(M)_{\sigma\circ\mathrm{rev},
        \tau\circ\mathrm{rev}};
        \notag
    \end{align}
    positive semidefiniteness is preserved under the simultaneous relabeling
    of rows and columns by the reflection.
\end{proof}

\section{Transfer maps}

\begin{definition}[MPO transfer map]\label{def:mpo_transfer_map}
    \lean{MPOTensor.transferMap}
    \leanok
    \uses{def:mpo_tensor}
    The \emph{transfer map} associated to $M$ is the linear map
    $\E_M:\MN{D}\to\MN{D}$ defined by
    \begin{align}
        \E_M(X)
        &=\sum_{i,j=0}^{d-1}M^{ij}X(M^{ij})^\dagger.
        \label{eq:mpdo_transfer_map}
    \end{align}
\end{definition}

\begin{lemma}[Identification with the doubled-index MPS transfer map]\label{lem:mpo_transfer_eq_mps}
    \lean{MPOTensor.transferMap_eq_toMPSTensor}
    \leanok
    \uses{def:mpo_to_mps, def:mpo_transfer_map, def:transfer_map}
    The transfer map of an MPO tensor agrees with the transfer map of its
    doubled-index MPS view: $\E_M=\E_{M^{\mathrm{MPS}}}$.
\end{lemma}

\begin{proof}
    \leanok
    \uses{def:mpo_to_mps, def:mpo_transfer_map, def:transfer_map}
    After identifying the single physical index of the doubled tensor with a
    pair $(i,j)$, for every $X\in\MN{D}$,
    \begin{align}
        \E_{M^{\mathrm{MPS}}}(X)
        &=
        \sum_{(i,j)}
        M^{\mathrm{MPS}}_{(i,j)}X
        (M^{\mathrm{MPS}}_{(i,j)})^\dagger
        =
        \sum_{i,j=0}^{d-1}M^{ij}X(M^{ij})^\dagger
        =
        \E_M(X).
        \notag
    \end{align}
\end{proof}

\begin{theorem}[Positivity of the MPO transfer map]\label{thm:mpo_transfer_pos}
    \lean{MPOTensor.transferMap_pos}
    \leanok
    \uses{def:mpo_transfer_map}
    If $X\ge0$, then $\E_M(X)\ge0$.
\end{theorem}

\begin{proof}
    \leanok
    \uses{lem:mpo_transfer_eq_mps}
    By Lemma~\ref{lem:mpo_transfer_eq_mps}, $\E_M$ is the transfer map of the
    doubled-index MPS tensor. The corresponding MPS transfer map is positive,
    so $\E_M(X)\ge0$ whenever $X\ge0$.
\end{proof}

\section{MPDOs and LPDOs}

\begin{definition}[MPDO]\label{def:mpdo}
    \lean{MPOTensor.IsMPDO}
    \leanok
    \uses{def:mpo_family}
    An MPO tensor is an \emph{MPDO} if, for every $N\geq1$,
    \begin{align}
        \rho^{(N)}(M)
        &\ge0.
        \label{eq:mpdo_positivity}
    \end{align}
    This is the global positivity condition of
    \cite[Section~4]{Cirac2016MPDO_arXiv}.
\end{definition}

\begin{definition}[LPDO]\label{def:lpdo}
    \lean{MPOTensor.IsLPDO}
    \leanok
    \uses{def:mpo_tensor}
    An MPO tensor is an \emph{LPDO} if there exist a Kraus dimension $d_K$,
    an inner bond dimension $D'$, a bond-space identification
    $e:\{0,\ldots,D-1\}\simeq
    \{0,\ldots,D'-1\}\times\{0,\ldots,D'-1\}$
    (equivalently, $D$ and $(D')^2$
    have the same cardinality), and matrices $A^{(i,k)}\in\MN{D'}$ such that,
    for all physical indices $i,j$,
    \begin{align}
        M^{ij}
        &=
        \left(\sum_{k=0}^{d_K-1}
            A^{(i,k)}\otimes\overline{A^{(j,k)}}\right)_{e,e},
        \label{eq:mpdo_lpdo_factorization}
    \end{align}
    where $\otimes$ denotes the Kronecker product, $\overline{(\cdot)}$
    denotes entrywise complex conjugation, and $(\cdot)_{e,e}$ denotes
    reindexing along $e$ back to $\{0,\ldots,D-1\}$.
    See \cite[Section~4.3]{Cirac2016MPDO_arXiv}.
    % Coefficient identity:
    % M^{ij}_{\alpha\beta}
    %   = \sum_{k,a,a',b,b'} e^{aa'}_{\alpha}
    %       A^{(i,k)}_{ab}\overline{A^{(j,k)}_{a'b'}}
    %       (e^{-1})^{\beta}_{bb'}.
    % Both sides have boundary signature (1,1,1,1): alpha and beta are the
    % west and east virtual indices, while i and j are the upper and lower
    % physical indices.  The west e wedge splits alpha into a,a'; the east
    % e^{-1} wedge combines b,b' into beta.  The unique inter-row pairleg is
    % the ancillary contraction k; a,a',b,b' are the remaining contracted
    % virtual indices.  No east cup is present, since beta remains free.
    The local contraction is
    \begin{tenkzequation}
        \tenkzkernel
        \begin{tenkz}[rows={wire}, cols=1, west=open, east=open]
            \tn[skin=mpo, ports={90:physical:$i$, 270:physical:$j$}]{M}
        \end{tenkz}
        $ = $
        \begin{tenkz}[rows={op,op}, cols=3]
            \tnfuse[at=(1,1), name=e, skin=pill]{e} \tnwire{open w}{e.180}
            \tn[at=(1,2), skin=box, ports={90:physical:$i$}]{A}
            \tn[at=(2,2), skin=box, ports={270:physical:$j$}]{\overline{A}}
            \tnfuse[at=(1,3), name=ei, skin=pill]{e^{-1}} \tnwire{ei.0}{open e}
        \end{tenkz}
    \end{tenkzequation}
\end{definition}

\begin{theorem}[LPDOs are MPDOs]\label{thm:lpdo_is_mpdo}
    \lean{MPOTensor.IsLPDO.isMPDO}
    \leanok
    \uses{def:lpdo, def:mpdo}
    Every LPDO tensor is an MPDO.
\end{theorem}

\begin{proof}
    \leanok
    \uses{def:lpdo, def:mpdo}
    By the mixed-product property
    $(A\otimes B)(C\otimes D)=(AC)\otimes(BD)$, the $N$-site product
    decomposes as a sum of Kronecker products after transporting along the
    bond-space identification $e$:
    \begin{align}
        M^{\sigma_1\tau_1}\cdots M^{\sigma_N\tau_N}
        &=
        \left(\sum_{\kappa}
            P_\kappa\otimes\overline{Q_\kappa}\right)\big|_e,
        \notag
    \end{align}
    where
    $P_\kappa
    =A^{(\sigma_1,\kappa_1)}\cdots A^{(\sigma_N,\kappa_N)}$
    and
    $Q_\kappa
    =A^{(\tau_1,\kappa_1)}\cdots A^{(\tau_N,\kappa_N)}$.
    Since the trace is invariant under the reindexing by~$e$, we apply
    $\tr(X\otimes Y)=\tr(X)\tr(Y)$ to obtain
    \begin{align}
        \rho^{(N)}(M)_{\sigma,\tau}
        &=
        \sum_\kappa
        \tr(P_\kappa)\overline{\tr(Q_\kappa)}
        =
        \sum_\kappa
        \psi_\kappa(\sigma)\overline{\psi_\kappa(\tau)},
        \notag
    \end{align}
    where $\psi_\kappa(\sigma)=\tr(P_\kappa)$.
    Hence
    $\rho^{(N)}(M)
    =\sum_\kappa\ketbra{\psi_\kappa}{\psi_\kappa}\ge0$.
\end{proof}
